\section*{Chapter \ref{ch:g10:reaction-progress-table} --- The Reaction-Progress Table}

\begin{solution}{exo:g10:reaction-progress-table:1}
\ce{H2}: $4.0 - 2x$; \ce{O2}: $3.0 - x$; \ce{H2O}: $2x$. Dihydrogen runs
out at $x = 2.0$, dioxygen at $x = 3.0$: $x_{\max} = \qty{2.0}{mol}$,
dihydrogen is limiting. \omterm{def:g10:reaction-progress-table:system}{Final state}: \qty{0}{mol} \ce{H2},
\qty{1.0}{mol} \ce{O2}, \qty{4.0}{mol} \ce{H2O}.
\end{solution}

\begin{solution}{exo:g10:reaction-progress-table:2}
Carbon would run out at $x = 3.0$, dioxygen at $x = 2.0$:
$x_{\max} = \qty{2.0}{mol}$, dioxygen is limiting. \omterm{def:g10:reaction-progress-table:system}{Final state}:
\qty{1.0}{mol} of carbon, no dioxygen, \qty{2.0}{mol} of carbon dioxide.
\end{solution}

\begin{solution}{exo:g10:reaction-progress-table:3}
$x_{\max} = \qty{0.30}{mol}$ (sulfur limiting). \omterm{def:g10:reaction-progress-table:system}{Final state}:
\qty{0.20}{mol} of iron, no sulfur, \qty{0.30}{mol} of iron sulfide.
\end{solution}

\begin{solution}{exo:g10:reaction-progress-table:4}
Dinitrogen runs out at $x = 1.0$, dihydrogen at $x = 2.4/3 = 0.80$:
$x_{\max} = \qty{0.80}{mol}$. \omterm{def:g10:reaction-progress-table:system}{Final state}: \qty{0.20}{mol} \ce{N2}, no
\ce{H2}, $2 \times 0.80 = \qty{1.6}{mol}$ \ce{NH3}.
\end{solution}

\begin{solution}{exo:g10:reaction-progress-table:5}
(b): $4/2 = 2/1$. In (a) carbon monoxide is limiting, in (c) too.
\end{solution}

\begin{solution}{exo:g10:reaction-progress-table:6}
$n(\ce{CH4}) = 32 / 16.0 = \qty{2.0}{mol}$, so $x_{\max} = \qty{2.0}{mol}$
for a \omterm{def:g10:reaction-progress-table:stoichiometric}{stoichiometric mixture}. Dioxygen: $2 \times 2.0 = \qty{4.0}{mol}$,
that is $4.0 \times 32.0 = \qty{128}{g}$. Water:
$\qty{4.0}{mol} \times \qty{18.0}{g/mol} = \qty{72}{g}$.
\end{solution}

\begin{solution}{exo:g10:reaction-progress-table:7}
$n(\ce{C3H8}) = 1000 / 44.0 = \qty{22.7}{mol}$; carbon dioxide
$3 \times 22.7 = \qty{68.2}{mol}$; $V = 68.2 \times 24.1 =
\qty{1.64e3}{L}$, about \qty{1.6}{m^3}.
\end{solution}

\begin{solution}{exo:g10:reaction-progress-table:8}
At $x = \qty{1.0}{mol}$: \qty{1.0}{mol} \ce{CH4}, \qty{1.0}{mol}
\ce{O2}, \qty{1.0}{mol} \ce{CO2}, \qty{2.0}{mol} \ce{H2O}. The dioxygen
is used up at $x = \qty{1.5}{mol}$: the reaction stops there, so larger
extents are never reached.
\end{solution}

\begin{solution}{exo:g10:reaction-progress-table:9}
$n_0(\ce{Zn}) = 1.30 / 65.4 = \qty{0.0199}{mol}$;
$n_0(\ce{Cu^{2+}}) = 0.10 \times 0.1000 = \qty{0.0100}{mol}$. The copper
\omterm{def:g9:ions:ion}{ions} are limiting: $x_{\max} = \qty{0.0100}{mol}$. Copper formed:
$0.0100 \times 63.5 = \qty{0.635}{g}$; zinc left:
$(0.0199 - 0.0100) \times 65.4 = \qty{0.65}{g}$. No copper \omterm{def:g9:ions:ion}{ions} remain:
the blue colour has gone.
\end{solution}

\begin{solution}{exo:g10:reaction-progress-table:10}
$n_0(\ce{Mg}) = 0.48 / 24.3 = \qty{0.0198}{mol}$, which would run out at
$x = 0.0198$; $n_0(\ce{H+}) = \qty{0.020}{mol}$, which would run out at
$x = 0.010$. The acid is limiting: $x_{\max} = \qty{0.010}{mol}$, so
\qty{0.010}{mol} of dihydrogen, $0.010 \times 24.1 = \qty{0.24}{L}$.
\end{solution}

\begin{solution}{exo:g10:reaction-progress-table:11}
A \omterm{def:g10:reaction-progress-table:stoichiometric}{stoichiometric mixture} with \qty{0.30}{mol} of sulfur needs
\qty{0.30}{mol} of iron; there is \qty{0.20}{mol} more, that is
$0.20 / 0.30 \approx \qty{67}{\percent}$ in excess.
\end{solution}

\begin{solution}{exo:g10:reaction-progress-table:12}
Equal amounts: $m(\ce{Fe})/55.8 = m(\ce{S})/32.1$, with
$m(\ce{Fe}) + m(\ce{S}) = \qty{10.0}{g}$. So
$m(\ce{Fe}) = 10.0 \times 55.8 / (55.8 + 32.1) = \qty{6.35}{g}$ and
$m(\ce{S}) = \qty{3.65}{g}$.
\end{solution}

\begin{solution}{exo:g10:reaction-progress-table:13}
\begin{enumerate}
  \item $n_0(\ce{CaCO3}) = 10.0 / 100.1 = \qty{0.0999}{mol}$, the only
        \omterm{def:g7:chemical-reactions:reactant}{reactant}: $x_{\max} = \qty{0.0999}{mol}$; it gives
        \qty{0.0999}{mol} of \ce{CaO} and of \ce{CO2}, that is
        $0.0999 \times 24.1 = \qty{2.41}{L}$ of carbon dioxide.
  \item $n_0(\ce{CaO}) = \qty{0.0999}{mol}$, $n_0(\ce{H2O}) = 1.8 / 18.0
        = \qty{0.100}{mol}$: the quicklime is (just) limiting; the
        \omterm{def:g6:pure-substances-and-mixtures:mixture}{mixture} is almost stoichiometric. \ce{Ca(OH)2}:
        $M = 40.1 + 2 \times (16.0 + 1.0) = \qty{74.1}{g/mol}$, mass
        $0.0999 \times 74.1 = \qty{7.40}{g}$.
\end{enumerate}
\end{solution}

\begin{solution}{exo:g10:reaction-progress-table:14}
$n_0(\ce{H+}) = \qty{0.050}{mol}$ in each flask; it would run out at
$x = \qty{0.025}{mol}$. Magnesium: 0.15, 0.30, 0.45, 0.60, 0.75,
\qty{0.90}{g} give 0.0062, 0.0123, 0.0185, 0.0247, 0.0309,
\qty{0.0370}{mol}. In the first four flasks magnesium is limiting
and the dihydrogen is $n(\ce{Mg}) \times 24.1$: 0.15, 0.30, 0.45 and
\qty{0.60}{L}. In the last two the acid is limiting: $0.025 \times 24.1
= \qty{0.60}{L}$. The plot is a straight line through the origin up to
about \qty{0.61}{g} of magnesium, then a horizontal plateau at
\qty{0.60}{L}.
\end{solution}

\begin{solution}{exo:g10:reaction-progress-table:15}
$n_0(\text{ethanol}) = 4.6 / 46.0 = \qty{0.100}{mol}$ (would run out at
$x = 0.100$); $n_0(\ce{O2}) = 4.8 / 24.1 = \qty{0.199}{mol}$ (would run
out at $x = 0.199/3 = 0.0664$). Dioxygen is limiting,
$x_{\max} = \qty{0.0664}{mol}$. Ethanol left: $0.100 - 0.0664 =
\qty{0.0336}{mol}$, that is $0.0336 \times 46.0 = \qty{1.5}{g}$. Carbon
dioxide: $2 \times 0.0664 = \qty{0.133}{mol}$, that is
$0.133 \times 24.1 = \qty{3.2}{L}$.
\end{solution}

\begin{solution}{pb:g10:reaction-progress-table:1}
\textbf{1.} Left: 2 Na, 6 N. Right: 2 Na, $3 \times 2 = 6$ N. Balanced.

\textbf{2.} Left: 10 Na, 2 K, 2 N, 6 O. Right: K: 2; Na: $5 \times 2 =
10$; O: $1 + 5 = 6$; N: 2. Balanced.

\textbf{3.} The periodic-table chapter (\cref{ch:g9:periodic-table-first-look}):
sodium reacts violently with water, giving a corrosive solution and
dihydrogen. After the crash the bag is in contact with the skin, the
eyes and the moist \omterm{def:g4:air-a-mixture-of-gases:air}{air} of breath.

\textbf{4.} $M(\ce{NaN3}) = 23.0 + 3 \times 14.0 = \qty{65.0}{g/mol}$;
$M(\ce{KNO3}) = 39.1 + 14.0 + 3 \times 16.0 = \qty{101.1}{g/mol}$.

\textbf{5.} \ce{NaN3}: $1.00 - 2x$; \ce{Na}: $2x$; \ce{N2}: $3x$.

\textbf{6.} $1.00 - 2x = 0$ at $x_{\max} = \qty{0.500}{mol}$; sodium
azide, the only \omterm{def:g7:chemical-reactions:reactant}{reactant}, is limiting.

\textbf{7.} \ce{Na}: $2 \times 0.500 = \qty{1.00}{mol}$; \ce{N2}:
$3 \times 0.500 = \qty{1.50}{mol}$.

\textbf{8.} $1.50 \times 24.1 = \qty{36.2}{L}$.

\textbf{9.} \ce{Na}: $1.00 - 10x$; \ce{KNO3}: $n - 2x$; \ce{K2O}: $x$;
\ce{Na2O}: $5x$; \ce{N2}: $x$.

\textbf{10.} $1.00/10 = n/2$, so $n = \qty{0.200}{mol}$, that is
$0.200 \times 101.1 = \qty{20.2}{g}$.

\textbf{11.} $x_{\max} = \qty{0.100}{mol}$: no sodium, no potassium
nitrate, \qty{0.100}{mol} \ce{K2O}, \qty{0.500}{mol} \ce{Na2O},
\qty{0.100}{mol} \ce{N2}.

\textbf{12.} Sodium would run out at $x = 0.100$, potassium nitrate at
$x = 0.150/2 = 0.075$: the nitrate is limiting, $x_{\max} =
\qty{0.075}{mol}$, and $1.00 - 10 \times 0.075 = \qty{0.25}{mol}$ of
sodium \omterm{def:g9:periodic-table-first-look:metal}{metal} would be left in the bag, ready to react with moisture.

\textbf{13.} $1.50 + 0.100 = \qty{1.60}{mol}$ of nitrogen per \omterm{def:g10:the-mole:amount}{mole} of
sodium azide.

\textbf{14.} $n = 60 / 24.1 = \qty{2.49}{mol}$.

\textbf{15.} $2.49 / 1.60 = \qty{1.56}{mol}$ of sodium azide.

\textbf{16.} $1.56 \times 65.0 = \qty{101}{g}$.

\textbf{17.} $0.200 \times 1.56 = \qty{0.311}{mol}$, that is
$0.311 \times 101.1 = \qty{31.5}{g}$ of potassium nitrate.

\textbf{18.} $2.49 / 1.50 = \qty{1.66}{mol}$, that is
$1.66 \times 65.0 = \qty{108}{g}$: the second reaction saves about
\qty{7}{g} of sodium azide.

\textbf{19.} A \omterm{def:g10:reaction-progress-table:limiting}{total reaction} stops only when its \omterm{def:g10:reaction-progress-table:limiting}{limiting reactant} is
used up; here sodium azide is the only \omterm{def:g7:chemical-reactions:reactant}{reactant}, so in the \omterm{def:g10:reaction-progress-table:system}{final state}
its amount is $1.00 - 2x_{\max} = 0$. Nothing toxic is left in the bag.

\textbf{20.} About \qty{101}{g} of sodium azide fill a \qty{60}{L}
airbag.
\end{solution}
